
\documentclass[11pt]{article}
\usepackage[a4paper,margin=0.92in]{geometry}
\usepackage[T1]{fontenc}
\usepackage[utf8]{inputenc}
\usepackage{lmodern,microtype,amsmath,amssymb,amsthm,graphicx,booktabs,float,hyperref,xcolor,enumitem}
\hypersetup{colorlinks=true,linkcolor=blue!45!black,citecolor=blue!45!black,urlcolor=blue!45!black}
\newtheorem{definition}{Definition}[section]
\newtheorem{proposition}[definition]{Proposition}
\newtheorem{theorem}[definition]{Theorem}
\newtheorem{remark}[definition]{Remark}
\newcommand{\C}{\mathbb C}
\newcommand{\PP}{\mathbb P}
\newcommand{\GG}{G}
\title{\textbf{Variational Principles for Finite-Slope Pandrosion Flows}\\[2mm]\large Implicit Descent, Local Acceleration, and Lyapunov Structure}
\author{Ivan Besevic\\Independent Mathematical Research}
\date{May 2026}
\begin{document}
\maketitle

\begin{abstract}
This paper proposes a variational interpretation of finite-slope Pandrosion flows.  Given an energy or residual functional $E$, the Pandrosion update is viewed as an implicit descent step driven by a probe-induced finite-slope surrogate.  The variational viewpoint explains continuous gates, local acceleration, and conservative behavior far from the root.  It also provides a language for proving energy decrease under structural assumptions on the probes, the local finite-slope matrix, and the line-search layer.
\end{abstract}

\section{Finite-slope flow viewpoint}
The Pandrosion correction
\[
        Q_G(a,b)\delta=-G(a)
\]
can be interpreted as a discrete flow direction.  Rather than approximating a classical derivative, $Q_G(a,b)$ encodes the geometry of a finite path from $a$ to $b$.  The update
\[
        x_{k+1}=x_k+\lambda_k\delta_k
\]
therefore behaves like a sampled flow whose vector field depends on the current probe geometry.

\begin{figure}[H]
\centering
\includegraphics[width=.96\textwidth]{figures/main_variational.png}
\caption{Variational interpretation of finite-slope Pandrosion: probes generate a tangent-like surrogate direction, the line search enforces descent, and the update follows a finite-slope flow toward a root or minimizer.}
\end{figure}

\section{Energy model}
Let $E:\C^n\to\mathbb R_+$ be an energy, for example
\[
        E(x)=\tfrac12\|G(x)\|^2+\rho\mathcal L(x)+\sigma\Pi(x),
\]
where $\mathcal L$ is a logarithmic scale energy and $\Pi$ penalizes equal-value or singular probes.  A Pandrosion direction is acceptable when it produces an energy decrease after line search:
\[
        E(x_k+\lambda\delta_k)\le E(x_k).
\]

\begin{definition}[Finite-slope descent direction]
A direction $\delta$ is a finite-slope descent direction at $x$ if it solves $Q_G(x,b)\delta=-G(x)$ for some admissible probe $b$ and if there exists $\lambda>0$ such that $E(x+\lambda\delta)<E(x)$.
\end{definition}

\section{Local descent principle}
\begin{proposition}[Implicit descent condition]
Assume $E$ is locally Lipschitz and that, in a neighborhood of a regular root, the finite-slope matrix $Q_G(x,b)$ is uniformly nonsingular and close to a stable local linearization in the finite-slope sense.  Then for $x$ sufficiently close to the root, the Pandrosion direction admits a positive step length $\lambda$ for which
\[
        E(x+\lambda\delta)<E(x).
\]
\end{proposition}
\begin{proof}
The finite-slope model gives $G(x)+Q_G(x,b)\delta=0$.  Uniform nonsingularity and finite-slope consistency imply that the residual decreases to first order along sufficiently small multiples of $\delta$.  Since $E$ is dominated locally by the residual term up to lower-order scale and penalty terms, a sufficiently small positive step decreases $E$.
\end{proof}

\begin{figure}[H]
\centering
\includegraphics[width=.78\textwidth]{figures/variational_landscape.png}
\caption{A synthetic energy landscape with a finite-slope trajectory.  The path is not a classical gradient method; its local direction is generated from finite-slope probes.}
\end{figure}

\section{Gated local acceleration}
The finite-slope Halley correction is best interpreted as a local variational acceleration.  Let $\omega(x)$ be a continuous gate depending on residual, condition, and logarithmic scale.  The accelerated direction is
\[
        \delta_{\omega}=(1-\omega)\delta_1+\omega(\delta_1+\delta_2),
\]
where
\[
        Q_G(a,b)\delta_2=-\tfrac12 H_G(\delta_1,\delta_1)
\]
with $H_G$ estimated by symmetric finite-slope probes.  Far from the root $\omega\approx 0$; near a stable root $\omega\approx 1$.

\begin{figure}[H]
\centering
\includegraphics[width=.78\textwidth]{figures/variational_gate.png}
\caption{Continuous acceleration gate.  High-order corrections are suppressed when the residual is large or the finite-slope matrix is poorly conditioned.}
\end{figure}

\section{Global Lyapunov picture}
A global variational theory should use an energy of the form
\[
        \mathcal V(x,i)=E_i(x)+\eta D_i(x),
\]
where $i$ is the active chart and $D_i$ is the chart distortion score.  Descent can then occur either by changing $x$ inside a chart or by switching charts while reducing $D_i$.

\begin{theorem}[Programmatic Lyapunov principle]
Suppose the line-search layer enforces $E_i(x_{k+1})\le E_i(x_k)$ within a chart and the atlas layer only accepts chart transitions that decrease $D_i$ up to a bounded tolerance.  Then $\mathcal V(x_k,i_k)$ is nonincreasing up to the accepted tolerance.
\end{theorem}

\section{Conclusion}
The variational viewpoint provides a clean mathematical bridge between raw Pandrosion, gated finite-slope Halley, projective scaling, and dynamic charts.  It also suggests a rigorous future path: prove convergence by constructing explicit Lyapunov energies adapted to finite-slope geometry.


\section{Surrogate action functional}
The finite-slope direction may be characterized by an implicit surrogate action.  Let
\[
        M_k(\delta)=\frac12\|G(x_k)+Q_G(x_k,b_k)\delta\|^2+\frac{\eta}{2}\|\delta\|^2.
\]
The raw Pandrosion correction is the zero-residual minimizer of the first term when $Q_G$ is nonsingular.  The regularized problem
\[
        \delta_k=\arg\min_\delta M_k(\delta)
\]
provides a variational interpretation even when $Q_G$ is ill-conditioned.  In the unregularized case this reduces to the equation $Q_G\delta=-G$.

\begin{proposition}[Regularized finite-slope direction]
For $\eta>0$, the surrogate action admits the normal equation
\[
        (Q_G^*Q_G+\eta I)\delta=-Q_G^*G(x_k).
\]
When $Q_G$ is square and nonsingular and $\eta\to 0$, the solution converges to the Pandrosion correction.
\end{proposition}

The practical engines discussed in the Pandrosion program do not need to explicitly solve this normal equation.  The proposition is conceptual: it shows that the finite-slope correction can be embedded into an energy-minimization framework.

\section{Energy decrease with line search}
The line-search layer is the operational mechanism that enforces the variational principle.  The algorithm tries a finite grid $\Lambda$ and chooses
\[
        \lambda_k\in\arg\min_{\lambda\in\Lambda} E(x_k+\lambda\delta_k).
\]
This makes descent robust even when the finite-slope model is only approximate.  In this sense, the Pandrosion correction supplies a direction and the line-search layer supplies the variational certificate.

\begin{theorem}[Discrete descent under accepted steps]
Suppose $\Lambda$ contains $0$ and the algorithm always accepts the grid minimizer of $E(x_k+\lambda\delta_k)$.  Then the energy sequence is nonincreasing.  If, in addition, the grid contains a positive step that strictly decreases energy whenever $x_k$ is not stationary, then the method cannot remain indefinitely outside the stationary set without decreasing $E$.
\end{theorem}
\begin{proof}
Since $0\in\Lambda$, the minimizer over $\Lambda$ cannot have energy larger than $E(x_k)$.  If a strict descent step exists, the minimizer has strictly smaller energy.  Iterating gives monotonicity and excludes stagnation away from the stationary set under the stated assumption.
\end{proof}

\section{Finite-slope curvature and second variation}
The finite-slope Halley layer estimates a second variation without constructing an analytic Hessian.  Along a direction $\delta$ one estimates
\[
        H_G(\delta,\delta)\approx G(x+h\delta)-2G(x)+G(x-h\delta).
\]
The correction
\[
        Q_G\delta_2=-\frac12 H_G(\delta_1,\delta_1)
\]
then cancels the leading second-order residual.  From the variational viewpoint, this is a local improvement of the surrogate action.  It should only be trusted when the first-order finite-slope model is already stable.

\section{Noise, probes, and robustness}
Probe geometry introduces variability.  Different probes generate different surrogates and therefore different descent directions.  This is not a defect; it is a source of geometric information.  A robust variational Pandrosion engine treats probes as a family of local models and selects the model that minimizes a combination of residual decrease, conditioning, curvature, and chart distortion.

\section{Connection with continuous flows}
If the step length becomes small and probes vary smoothly with $x$, the discrete update suggests a continuous flow
\[
        \dot x=-V(x),\qquad V(x)=-Q_G(x,b(x))^{-1}G(x).
\]
This is not a classical gradient flow unless $Q_G$ is symmetric positive and aligned with a metric.  However, it is a natural finite-slope flow.  The variational problem is to identify a metric or energy under which $V$ is descent-like.

\section{Implementation guidance}
The variational interpretation suggests three practical rules:
\begin{enumerate}[leftmargin=2em]
\item use conservative gates far from a root;
\item line-search every high-order correction;
\item treat chart and probe penalties as part of the energy, not as external heuristics.
\end{enumerate}
These rules explain why the more stable engines in the Pandrosion sequence separate global geometric control from local acceleration.

\section{Open problems}
The main open problem is to construct explicit energies $E$ for broad classes of polynomial systems such that finite-slope Pandrosion directions are provably descent directions in large regions.  A second problem is to characterize when the gated finite-slope Halley correction preserves descent.  A third is to connect the dynamic atlas score with a Lyapunov function on an extended state space.

\begin{thebibliography}{9}
\bibitem{traub} J. F. Traub, \emph{Iterative Methods for the Solution of Equations}, Prentice-Hall, 1964.
\bibitem{householder} A. S. Householder, \emph{The Numerical Treatment of a Single Nonlinear Equation}, McGraw-Hill, 1970.
\bibitem{doCarmo} M. do Carmo, \emph{Riemannian Geometry}, Birkhauser, 1992.
\bibitem{amari} S. Amari, \emph{Information Geometry and Its Applications}, Springer, 2016.
\end{thebibliography}
\end{document}
